\section{限制、工具与课堂 Q\&A：把研究原型接到真实系统}

课程最后没有用 SOTA 数字收尾，而是集中讨论哪些接口尚未解决：短 horizon、单时间尺度、toy environments、image goal、真实机器人、agent architecture 与 policy distillation。Q\&A 也反复提醒，JEPA 是训练框架而非新网络类型，world model 可以与 diffusion、Transformer、RL policy 和 classical control 组合。

\subsection{四个公开限制对应四条研究路线}

slide 53 把 limitation 写成 research opportunity。短 horizon 要求更稳的 rollout 或 closed-loop correction；单一时间尺度要求 hierarchical state；toy setup 要求真实感知与 domain shift；不现实的 goal image 要求语言、任务约束或 reward interface。任何一条都不仅是``扩大模型''即可解决。

\lecturefigure{slide-053.jpg}{LeWorldModel 的限制：短 horizon、单尺度、toy environment 与 goal interface}{CS25 V6 official deck, p.~53}

\begin{knowledgebox}{从 limitation 到系统模块}
\begin{tabularx}{\textwidth}{L{0.24\textwidth} X}
\toprule
限制 & 需要补上的系统能力 \\
\midrule
short-term planning & uncertainty-aware rollout、replanning、policy/world-model hybrid \\
single temporal level & hierarchical latent、event abstraction、long-term memory \\
toy environments & robust perception、sim-to-real、partial observability、safety constraints \\
unrealistic image goal & language/task encoder、reward model、constraint specification \\
\bottomrule
\end{tabularx}
\end{knowledgebox}

\teachervoice{00:55:14--00:57:15，Lucas 明确把这些项目称为 limitations or research opportunities：当前 planner 主要是短期、单层 temporal representation，环境和目标接口仍远离开放世界机器人。}

\subsection{stable-worldmodel：可复现比较也是方法贡献}

slide 54 展示一个统一库，将 LeWM、PLDM、DINO-WM 等 baselines，与 CEM、MPPI、Adam 等 solvers，以及 Push-T、TwoRoom、OG-Bench、DMC 等 environments 放到同一接口。它的科学价值是减少每篇论文因不同数据预处理和 planner 实现造成的不可比性。

\lecturefigure{slide-054.jpg}{stable-worldmodel：统一 baseline、solver 与 environment 的实验层}{CS25 V6 official deck, p.~54}

\teachervoice{00:57:15--00:59:16，老师邀请学生使用和贡献 `stable-worldmodel`，并强调 baselines、solvers、environments、测试与文档都已统一。这里的重点是 reproducibility，而非把某个库当作唯一标准。}

\begin{importantbox}{统一 API 不能替代统一协议}
同一库仍需固定 observation resolution、action normalization、dataset split、planner budget、random seeds 与 success definition。API 解决代码接口，benchmark protocol 才解决科学可比性；两者缺一不可。
\end{importantbox}

\subsection{Physical AI：为什么动作后果预测是独立训练目标}

Q\&A 中 Lucas 对当前 VLA 是否具有可靠 world understanding 持怀疑态度：若模型没有被训练去预测自身动作后果，它可能在熟悉模仿任务中表现良好，却缺少面向新交互的显式检查。Hazel 补充 world model 还可作为 policy evaluator。这个观点应作为 speaker judgment，而非对所有 VLA 的实证定论。

\teachervoice{00:59:16--01:01:17，Lucas 认为 physical agent 若不能预测 action consequence，就难以可靠互动；Hazel 补充 world model 可以评估候选 policy。讲义提醒：这是前瞻判断，需要真实机器人对照实验验证。}

\begin{warningbox}{VLA 与 world model 不是互斥类别}
VLA 可以内部包含 predictive objective、latent simulator 或 value/world model；world model 也可用语言与视觉 Transformer 实现。真正差别在训练接口：模型是否显式学习 $p(s_{t+1}\mid s_t,a_t)$，以及部署时是否用该预测检查或搜索动作。
\end{warningbox}

\subsection{Masking 是否必要：必要的是排除 shortcut}

学生问 prediction loss 是否不足。Hazel 的回答很谨慎：对所有 world model，masking 不是逻辑上唯一方法；但在 object-centric setup 中，无遮挡 prediction 容易学 self-dynamics，无法保证 interaction-based dynamics。masking 是一种把 shortcut 移出可行解空间的手段，而非普适定理。

\teachervoice{01:01:17--01:03:15，老师回答 ``masking necessary?''：prediction loss 本身可以学习某些 world model，但可能只学对象自己的运动；object masking 强化了 interaction dynamics，不应被写成所有系统必需的唯一 recipe。}

\begin{importantbox}{设计自监督目标的通用问题}
不要只问 ``loss 能否下降''，而要问：\textbf{最便宜的 shortcut 是什么？} 若目标任务需要关系、长期记忆或动作后果，训练时就要系统性移除只靠局部统计、身份泄漏或静态相关性即可完成的路径。
\end{importantbox}

\subsection{JEPA-native agent 与 LLM tool agent 的接口差异}

Lucas 把 JEPA-native agent 描述为 action-conditioned predictive model 加 classical control：输入当前 state 与动作，预测未来，再用 MPC 零样本搜索。LLM agent 常通过 text/tool-call loop 学习使用工具，未必显式学习连续物理动作的状态转移。两者可组合，例如语言模型分解目标，world model 负责低层 action consequence。

\teachervoice{01:03:19--01:05:20，课堂强调 JEPA-native agent 的不同点是 action-conditioned future model，而不是使用某种新 backbone。学到 dynamics 后，可以直接接 MPC，无需先把每个目标都 post-train 成 policy。}

\begin{knowledgebox}{Agent stack 的可组合分层}
语言模型可负责任务解释与高层 subgoal；world model 负责预测连续环境状态；planner/solver 负责搜索动作；policy 负责快速执行；verifier 与 safety shield 负责约束。把这些层分开，才能定位失败来自语义、动力学、优化还是执行。
\end{knowledgebox}

\subsection{JEPA 不是 Transformer 的反义词}

``JEPA 是否比 Transformer 更少 hallucination'' 这个问题混合了 framework 与 architecture。Lucas 明确指出 LeWorldModel 的 encoder 和 predictor 本身就是 Transformers；JEPA 描述训练时预测什么，Transformer 描述 token 如何交互。world model 仍会因数据不足、容量不足、模型误差或 controller 搜索到越界动作而产生错误未来。

\teachervoice{01:05:20--01:07:22，Lucas 纠正类别错误：JEPA 不是新 architecture，而是学习 world model 的 framework。即便使用 JEPA，坏数据或坏模型仍会 hallucinate；只是错误类型与 LLM 的无 grounding 文本生成并不完全相同。}

\begin{warningbox}{World-model hallucination 的四个来源}
\begin{enumerate}
  \item encoder 丢失决策相关状态；
  \item predictor 在长 horizon 累积误差；
  \item training distribution 缺少关键 interaction；
  \item planner 提出环境无效或越界动作。
\end{enumerate}
因此 ``更 grounded'' 是可检验假设，不是由 JEPA 名称自动保证的属性。
\end{warningbox}

\subsection{Diffusion、MPC 与 policy distillation 可以共存}

另一个问题把 diffusion model 与 world model 当成竞争品。讲者指出 diffusion 可以作为 world-model predictor、action proposal 或 policy；world model 是功能角色，不限定概率模型家族。对于长 horizon，系统还可以用昂贵 MPC 生成高质量行为，再蒸馏成快速 policy，熟悉场景走反射路径，困难场景回退到规划。

\teachervoice{01:07:22--01:10:59，课堂给出两条系统答案：diffusion 可以嵌入 JEPA/world-model pipeline；未来 agent 可能像 fast/slow system，简单情况直接 policy，困难情况调用 MPC，并把规划经验再蒸馏回 policy。}

\begin{importantbox}{Fast policy 与 slow planner 的工程闭环}
离线 world model 学 dynamics；在线 MPC 在新目标上产生动作；成功轨迹进入 replay buffer；policy distillation 学习快速映射；低置信度或分布外状态重新触发 planner。这个闭环把模型、控制和持续数据收集连接起来，但也引入 model bias、policy drift 与安全验证问题。
\end{importantbox}

\subsection{本章小结}

真实系统不会在 JEPA、Transformer、diffusion、MPC 与 policy 中五选一。它们分别回答训练目标、网络结构、概率建模、在线搜索和快速执行。当前研究原型的主要缺口是 long horizon、hierarchy、真实环境与目标接口；统一 benchmark tooling 能提高比较质量，但不能替代真实部署与安全证据。

